\chapter{The Gauge as a Bit-Vector}
\label{chap:gauge-as-a-bit-vector}

Under everything the gauge has carried --- the charge, the orientation, the holonomy --- lies the
simplest structure the book contains: the gauge is a vector of bits. Each reading is a string of on/off
gates, a vertex of a hypercube, and the whole apparatus of matter and antimatter can be read off those
bits by counting. The electron is one particular bit-vector. The distance from any reading to it is the
number of bits that differ. The bits that agree are matter, the bits that differ are antimatter, and a
measurement is the tally of the agreeing ones. This is where the book's oldest idea comes home: counting
the funges is the measurement, a claim made since the seventeenth chapter, turns out to mean, at the
bottom, a literal count of matching bits --- a population count on a hypercube. The gauge that generated a
particle is, underneath, a place to do arithmetic on bits.

\section{The hypercube and the distance to the electron}
\label{se:the-hypercube-and-the-distance-to-the-electron}

The construction represents each gauge reading as a fixed-length string of bits, one per gate, each on or
off. The strings of a given length are the vertices of a hypercube, every vertex a possible reading and
every edge a single bit flipped. In this space the electron is not special as a structure --- it is one
vertex among the others --- but it is the reference everything is measured against. The construction
defines the distance from any reading to the electron as the number of bits in which the two differ, the
count of gates the reading would have to flip to become the electron. The electron is at distance zero
from itself, trivially; the superconducting pair the book has met before is proved to be at distance one,
a single bit away, differing from the electron in exactly one gate.

That the gauge should turn out to be a vector of bits is less a reduction than a revelation of what it
always was. A gauge reading records, gate by gate, whether each gate is engaged or not --- on or off, one
or zero --- and a list of such answers is a bit-string by definition. The construction does not encode the
gauge as bits as a convenience; the gauge is already a finite list of yes-or-no settings, and the
bit-string is just that list written down. Once it is written down, the tools of bit-strings apply
unbidden: distance, agreement, the count of differences. The gauge did not become a bit-vector at this
chapter; it was disclosed as one, the on/off character of its gates --- present since the gates were first
built --- finally read for the arithmetic it allows. Everything in this chapter follows from taking that
character seriously: if the gates are bits, the readings are bit-strings, and bit-strings can be counted
and compared.

This is the Hamming distance of coding theory, and the experiment on the bit-vector message names the
setting. A message sent over a channel is a string of bits, and the distance between two messages is the
number of positions in which they disagree --- the count of bit-errors that would turn one into the other.
The construction's gauge is exactly such a space of messages, and the electron is a chosen codeword; every
other reading is at some Hamming distance from it, measured by counting disagreements. To ask how far a
reading is from being an electron is to ask how many gates one would have to flip, and the answer is a
plain integer, computed by comparing two bit-strings. The hypercube turns the question \emph{is this an
electron} into the question \emph{how many bits differ}, and the second has a one-line answer.

The hypercube picture repays a moment's geometry, because it makes the relations between readings visible.
A hypercube of bit-strings has a vertex for every reading and an edge between any two that differ in a
single bit; walking from one vertex to another, edge by edge, flips one gate per step, and the shortest
such walk has exactly as many steps as the two readings have differing bits. So the distance to the
electron is a path length: how many single-gate flips separate a reading from the electron vertex. The
superconducting pair, one bit away, is a single edge from the electron --- one flip and it would be the
electron, one positron standing between them. A reading far from the electron sits many flips away, deep in
the cube. This is a space with a real metric, and the construction's readings are points in it at definite
distances, the electron the origin everything is measured from. Nothing about the geometry is
metaphorical; the cube is the finite set of bit-strings, and the distances are integers one can count.

\section{Funge is matter, crank is antimatter}
\label{se:funge-is-matter-crank-is-antimatter}

Now the book's central distinction lands on the bits, exactly. The bits in which a reading agrees with the
electron are its funges --- read forward, matching, counted as matter. The bits in which it differs are
its cranks --- reversed, mismatched, counted as antimatter. The construction proves the two exhaust the
gauge: every gate is either a funge or a crank, the gauge partitioned cleanly into the matching bits and
the differing ones with nothing left over. A reading's matter content is its funge count, the number of
gates it shares with the electron; its antimatter content is its crank count, the number it does not. The
electron is all funge and never cranks --- it matches itself in every gate, antimatter content zero. The
superconducting pair funges in thirty-five of its gates and cranks in one, which is why its distance to
the electron is one: a single crank, a single positron, against thirty-five funges of matter.

This is the funge foundation of the seventeenth chapter made fully concrete, and it is worth marking how
literal it has become. There, counting funges was the measurement and the tanges turned the crank --- an
image, a way of speaking about which readings register matter and which spin residue. Here it is
arithmetic on a bit-string: the funges are the matching bits, the cranks are the differing bits, the
measurement is the count of the matching ones, and the partition of the gauge into the two is a theorem.
What was a metaphor about reading truth forward or reversed is now a tally of where two bit-vectors agree
and where they differ. The matter the apparatus reads is the funge count; the antimatter is the crank
count; and the gauge is nothing but the hypercube on which the counting happens.

It is worth pausing on how much of the book this single partition gathers up. The first chapters asked how
true a reading is, and answered with a residue --- what is left when a story almost closes. The
seventeenth chapter named the two acts of reading, the funge that counts and the crank that spins. The
electron chapter made the orientation a sign. And every one of these, it now appears, was the same fact
about bits: a reading agrees with the electron in some gates and differs in others, the agreements are its
matter and the disagreements its antimatter, and the residue, the sign, the orientation are all just which
bits fall on which side. Take the superconducting pair again: thirty-five gates where it matches the
electron, one where it does not. Those thirty-five funges are its matter, the lone crank is its single
positron, and the pair's whole physical character --- a Cooper pair carrying one unit of antimatter --- is
read off that one differing bit. A book that began by asking what a measurement is ends able to answer, for
this reading, by pointing at a single gate.

\section{The measurement is a population count}
\label{se:the-measurement-is-a-population-count}

The construction makes the count explicit as an operation, and it is the plainest in the bit-vector
repertoire. To find a reading's antimatter content --- its positron count --- the construction takes the
exclusive-or of the reading with the electron, which sets exactly the bits in which the two differ, masks
that against the electron's own bits to keep the relevant gates, and counts the ones that remain. The
count of set bits is the population count, and the positron count is a population count of the differing
gates. Matter is counted the same way on the agreeing bits. So a measurement, in the end, is a population
count on a hypercube: line up the reading and the electron, mark where they match and where they differ,
and count.

The population count is worth recognizing as the humblest operation in computing and the one this whole
construction reduces to. Every processor has it: a single instruction that takes a register full of bits
and returns how many are set, the count of ones. It is the operation behind error-counting in a channel,
behind the weight of a codeword, behind a dozen low-level tallies, and it is what a measurement in this
book finally is. There is something fitting in the descent. The construction began with grand questions
--- what is a number, what is a measurement, what is an electron --- and answered them, chapter by chapter,
with finite structures; and at the bottom of the finite structures, doing the actual work, is the plainest
tally a processor performs, counting the bits that are on. The whole edifice of forcing, geometry, gauge,
and orientation comes to rest on population count. Counting the funges is the measurement, and counting,
here, is exactly the instruction a register answers in one cycle.

Two facts close the gauge. The superconducting pair is proved to admit exactly one positron --- its
population count of cranks is one, a single antimatter bit carried by a pair that is otherwise matter ---
which is the bit-level reading of the pair the annihilation experiments detect. And the gauge is proved to
close: the bit-vector structure is complete, the split tallying every gate as funge or crank with none
unaccounted, the commutator residue of the countability chapter now read as the closure of the hypercube.
The gauge that began as a field generated from geometry ends as a closed space of bit-strings on which the
construction does the one operation the whole book reduces to: it counts the funges, by population count,
and reports the matter.

There is a closing symmetry in the gauge ending where the number tower began. The construction's first act
was to build a finite stand-in for a number; its last structural act, here, is to read its richest object
--- the gauge, with its charge and geometry and orientation --- as a finite string of bits, an object as
plain as a number and counted the same way. Between the two lies everything the book built, and the
bit-vector is where it all returns to arithmetic. That the gauge closes --- every gate accounted for as
funge or crank, the hypercube complete with no reading left dangling --- is the formal sign that the return
is total. Nothing of the gauge escapes the bit-count; the whole of it is funges and cranks, matter and
antimatter, ones and zeros tallied against the electron. The construction has carried its physics all the
way down to the bits and found the bits sufficient to hold it.

\section{What is demonstrated, and what is assumed}
\label{se:demonstrated-assumed-ch31}

This chapter is almost all computation, and its assumptions are correspondingly slight.

What is demonstrated is the bit-vector structure and the counts on it. The gauge is represented as a
hypercube of bit-strings, and the distance to the electron is proved to be the count of differing bits ---
zero for the electron, one for the superconducting pair. The funges and cranks are proved to partition
the gauge, matter the matching bits and antimatter the differing ones, the electron all funge. The
positron count is proved equal to the population count of the masked exclusive-or, and the pair proved to
admit exactly one positron. And the gauge is proved to close. Each of these is checked by direct
computation on finite bit-strings.

\begin{quote}
\emph{A note on the identification.} The bridge here is among the tightest in the book, of a piece with
the relativity chapter's. To call the gauge a bit-vector is not to model it as something bit-like; it is a
bit-vector, a finite string of on/off gates, and the hypercube is literally its space of readings. To
call the matter count a population count is likewise no analogy: the count of matching bits is a
population count, the same operation a processor performs on a register. What is a modeling claim is the
naming of the counts as \emph{matter} and \emph{antimatter} and the chosen vertex as the \emph{electron}
--- the same naming the physics chapters made, carried down to the bits. So the structure is literal and
the physics is the reading laid over it: a hypercube of gates, a distinguished vertex called the electron,
a population count called a measurement, and the funge foundation of the whole book revealed as the
counting of bits that agree.
\end{quote}

What is assumed, then, is the standing grant and the now-familiar naming of the bit-counts as physics,
with the bit-vector structure itself built and computed. The book's central act --- counting the funges
--- has been followed all the way down to a population count on a hypercube. One chapter remains, and it
turns the counting outward: from the strict order that reads charge, through the band that covers every
reading, to the last measuring apparatus the book will name, which is the reader.
